% =============================================================
%  Chapter — Standards & Standardization Bodies
%  Key players in the standardization and publication landscape
%  relevant to HVDC: IEC (incl. document types IS/TS/TR/PAS), ISO, CEN/CENELEC, IEEE, CIGRE, EPRI, ESIG.
% =============================================================
\sectionslide{Standardization Bodies}{Who writes the standards and technical recommendations}

% ---- IEC ----------------------------------------------------
\begin{frame}{IEC\\International Electrotechnical Commission}
  \begin{block}{Role}
    The world's leading organization for international standards in electrical, electronic and related technologies.
  \end{block}
  \vspace{0.3em}
  \begin{itemize}
    \setlength{\itemsep}{0.4em}
    \item Founded 1906; headquartered in Geneva, Switzerland.
    \item Consensus-based, normative standards developed by national committees and technical committees (e.g. TC~115 for HVDC).
    \item Also runs conformity-assessment / certification schemes.
    \item IEC standards are normative for countries adopting the IEC standards.
  \end{itemize}
\end{frame}


% ---- IEC document types -------------------------------------
\begin{frame}{IEC Document Types\\IS, TS, TR and PAS}
  \noindent{\color{atGray}\small The prefix after ``IEC'' tells whether a document sets requirements or only gives guidance.}
  \vspace{0.6em}

  {\footnotesize
  \renewcommand{\arraystretch}{1.25}
  \begin{tabularx}{\linewidth}{@{}l l >{\raggedright\arraybackslash}X >{\raggedright\arraybackslash}X@{}}
    \toprule
    \multicolumn{1}{c}{\textbf{Type}} & \multicolumn{1}{c}{\textbf{Content}} & \textbf{Approval} & \textbf{Typical use / HVDC example} \\
    \midrule
    IEC (IS)   & Normative     & Full consensus: CDV + FDIS vote     & Ratings and tests; IEC~60700-1, IEC~62501 \\
    \atrowfill
    IEC TS     & Normative     & TC vote, 2/3 of P-members           & Emerging technology; IEC~TS~63291-1 \\
    \atrowhighlight
    IEC TR     & Informative   & TC vote, simple majority            & Guidance, state of the art; IEC~TR~60919-1 \\
    IEC PAS    & Interim       & TC vote, simple majority            & Fast-track industry or consortium work \\
    \bottomrule
  \end{tabularx}}

  \vspace{0.6em}
  \noindent{\color{atGray}\small Compliance can be claimed only with an IS or TS. A TR binds a vendor only where the specification writes its content in as requirements.}
\end{frame}


% ---- ISO ----------------------------------------------------
\begin{frame}{ISO\\International Organization for Standardization}
  \begin{block}{Role}
    International standards for all technical fields outside
    electrotechnology (which is covered by the IEC).
  \end{block}
  \vspace{0.3em}
  \begin{itemize}
    \setlength{\itemsep}{0.4em}
    \item Founded 1947; headquartered in Geneva, Switzerland.
    \item Members are national standards bodies, one per country (e.g.\ DIN, ANSI, BSI), roughly 170 in total.
    \item Covers management systems (ISO~9001, 14001, 45001), mechanical engineering, materials and testing; joint work with IEC (e.g.\ ISO/IEC~17025).
    \item Like IEC standards, ISO standards become normative only once adopted nationally or referenced in contracts.
  \end{itemize}
\end{frame}


% ---- CEN / CENELEC ------------------------------------------
\begin{frame}{CEN \& CENELEC\\European Committees for Standardization}
  \begin{block}{Abbreviations}
    {\footnotesize
    \begin{tabularx}{\linewidth}{@{}l@{\hspace{0.8em}}>{\raggedright\arraybackslash}X@{}}
      \textbf{CEN}     & Comit\'e Europ\'een de Normalisation --- European Committee for Standardization \\
      \textbf{CENELEC} & Comit\'e Europ\'een de Normalisation \'Electrotechnique --- European Committee for Electrotechnical Standardization \\
      \textbf{EN}      & Europ\"aische Norm / European Norm --- European Standard, published by CEN and CENELEC \\
    \end{tabularx}}
  \end{block}
  \vspace{0.3em}
  \begin{itemize}
    \setlength{\itemsep}{0.4em}
    \item Founded 1961 (CEN) and 1973 (CENELEC); Brussels; 34 national members.
    \item Every member \textbf{must} adopt an EN unchanged; conflicting standards are withdrawn.
    \item Most ENs adopt ISO work (via CEN) or IEC work (via CENELEC).
    \item Harmonized ENs give presumption of conformity with EU directives.
  \end{itemize}
\end{frame}


% ---- Designations -------------------------------------------
\begin{frame}{Standard Designations\\DIN EN, DIN EN ISO, DIN IEC}
  \noindent{\color{atGray}\small The prefix shows the origin of a standard and how widely it is binding; the number stays the same as in the original.}
  \vspace{0.6em}

  {\footnotesize
  \renewcommand{\arraystretch}{1.25}
  \begin{tabularx}{\linewidth}{@{}l >{\raggedright\arraybackslash}X >{\raggedright\arraybackslash}X l@{}}
    \toprule
    \textbf{Designation} & \textbf{Origin} & \textbf{Adoption path} & \textbf{Binding in} \\
    \midrule
    DIN               & German national          & DIN only                                   & Germany \\
    \atrowfill
    DIN EN            & European (CEN / CENELEC) & All CEN / CENELEC members                  & Europe \\
    DIN EN ISO        & International (ISO)      & ISO $\rightarrow$ CEN $\rightarrow$ DIN    & Worldwide + Europe \\
    \atrowfill
    DIN EN IEC / 6xxxx & International (IEC)     & IEC $\rightarrow$ CENELEC $\rightarrow$ DIN & Worldwide + Europe \\
    DIN ISO           & International (ISO)      & ISO $\rightarrow$ DIN, no EN               & Germany only \\
    \atrowhighlight
    DIN IEC           & International (IEC)      & IEC $\rightarrow$ DIN / DKE, no EN         & Germany only \\
    \bottomrule
  \end{tabularx}}

  \vspace{0.6em}
  \noindent{\color{atGray}\small Many HVDC documents are IEC~TS / TR that CENELEC has not adopted. For US projects, the IEC number itself is what counts.}
\end{frame}


% ---- IEEE ---------------------------------------------------
\begin{frame}{IEEE\\Institute of Electrical and Electronics Engineers}
  \begin{block}{Role}
    The world's largest technical professional organization; publisher of
    standards, journals and conference proceedings.
  \end{block}
  \vspace{0.3em}
  \begin{itemize}
    \setlength{\itemsep}{0.4em}
    \item Standards developed through the IEEE Standards Association (IEEE~SA).
    \item Especially influential in North American practice.
    \item Power \& Energy Society (PES) covers transmission, HVDC and FACTS.
    \item IEEE standards are not normative, unless adopted by individual network operators.
  \end{itemize}
\end{frame}


% ---- CIGRE --------------------------------------------------
\begin{frame}{CIGRE\\International Council on Large Electric Systems}
  \begin{block}{Role}
    A global collaborative community pooling operator, manufacturer and
    academic expertise on large electric power systems.
  \end{block}
  \vspace{0.3em}
  \begin{itemize}
    \setlength{\itemsep}{0.4em}
    \item Founded 1921; headquartered in Paris, France.
    \item Output: \textbf{Technical Brochures} (TBs) and conference publications.
    \item Work organized into \textbf{16 Study Committees} across four domains (HVDC sits in \textbf{SC~B4}).
    \item CIGRE Technical Brochures are authoritative guidance and state-of-the-art references.
    \item TBs frequently become the technical basis for later IEC/IEEE standards.
  \end{itemize}
\end{frame}


\begin{frame}{CIGRE\\Study Committees}
  \noindent{\color{atGray}\small 16 Study Committees grouped into four domains:}
  \vspace{0.6em}
  \begin{columns}[T,onlytextwidth]
    \begin{column}{0.48\textwidth}
      {\color{atAccent}\textbf{A — Equipment}}\\[0.35em]
      \footnotesize
      \textbf{A1} Rotating Electrical Machines\\[0.25em]
      \textbf{A2} Power Transformers \& Reactors\\[0.25em]
      \textbf{A3} Transmission \& Distribution Equipment
    \end{column}
    \begin{column}{0.48\textwidth}
      {\color{atAccent}\textbf{B — Subsystems}}\\[0.35em]
      \footnotesize
      \textbf{B1} Insulated Cables\\[0.25em]
      \textbf{B2} Overhead Lines\\[0.25em]
      \textbf{B3} Substations \& Electrical Installations\\[0.25em]
      {\color{atInk}\textbf{B4} DC Systems \& Power Electronics}\\[0.25em]
      \textbf{B5} Protection \& Automation
    \end{column}
  \end{columns}
  \vspace{0.8em}
\end{frame}


\begin{frame}{CIGRE\\Study Committees}
  \noindent{\color{atGray}\small 16 Study Committees grouped into four domains:}
  \vspace{0.6em}
  \begin{columns}[T,onlytextwidth]
    \begin{column}{0.48\textwidth}
      {\color{atAccent}\textbf{C — System}}\\[0.35em]
      \footnotesize
      \textbf{C1} System Development \& Economics\\[0.25em]
      \textbf{C2} System Operation \& Control\\[0.25em]
      \textbf{C3} System Environmental Performance\\[0.25em]
      \textbf{C4} System Technical Performance\\[0.25em]
      \textbf{C5} Electricity Markets \& Regulation\\[0.25em]
      \textbf{C6} Active Distribution \& DER
    \end{column}
    \begin{column}{0.48\textwidth}
      {\color{atAccent}\textbf{D — Materials}}\\[0.35em]
      \footnotesize
      \textbf{D1} Materials \& Emerging Test Techniques\\[0.25em]
      \textbf{D2} Information Systems \& Telecommunication
    \end{column}
  \end{columns}
  \vspace{0.8em}
\end{frame}


% ---- EPRI ---------------------------------------------------
\begin{frame}{EPRI\\Electric Power Research Institute}
  \begin{block}{Role}
    Independent, non-profit R\&D organization serving the electric power
    industry — research, not standards.
  \end{block}
  \vspace{0.3em}
  \begin{itemize}
    \setlength{\itemsep}{0.4em}
    \item Founded 1972; headquartered in Palo Alto, California.
    \item Funded by and for the industry; collaborates with 450+ companies across ~45 countries.
    \item Covers generation, delivery and end-use of electricity — reliability, efficiency, safety and environment.
    \item Output: research reports, technical guidelines and testing that inform utility practice.
  \end{itemize}
\end{frame}


% ---- ESIG ---------------------------------------------------
\begin{frame}{ESIG\\Energy Systems Integration Group}
  \begin{block}{Role}
    Non-profit technical \& educational association focused on integrating
    new resources and technologies into the power system.
  \end{block}
  \vspace{0.3em}
  \begin{itemize}
    \setlength{\itemsep}{0.4em}
    \item Formerly UVIG (Utility Variable-Generation Integration Group); renamed in 2018 after merging with iiESI.
    \item Focus: high levels of renewables / inverter-based resources, HVDC and emerging technologies on the grid.
    \item Members: utilities, operators, manufacturers, labs, universities and consultants.
    \item Output: reports, workshops, webinars and short courses — guidance and best practice, not normative.
  \end{itemize}
\end{frame}

